\section{Data and Measurement}\label{sec:data}

\subsection{Samples}

The main data are the PSID. Its core sample combines the Survey Research Center (SRC) sample, a national
probability sample of 1968 households, and the Survey of Economic Opportunity (SEO) sample, an oversample of low-income
households drawn in 1968 mostly in cities and in the South. The SEO is not a random sample of Black families: it is a
sample of low-income ones, and it provides most of the Black children (217 of the Black sons in the analysis come from
the SRC). Children are linked to the household in which they grew up and followed as adults, as heads or spouses of
their own households, from 25 to 55. Race is the race of the parental household head. Parental income is the mean of
the parental household's family income over the income years in which the child is 13--19, with at least three years
\citep[as in][]{lee2009trends, palomino2018one}; for ranks it is the mean at 12--17. The children's outcomes are their
own labor earnings and the family income of their own household. The benchmark sample has 1,995 Black and 2,758 White
sons (17,710 and 31,182 child-years at 25--65) and 2,367 Black and 2,852 White daughters.\footnote{Yearly amounts below
\$100 or above \$150,000 in 1967 dollars are dropped in both generations \citep{palomino2018one}. Code:
\texttt{6.IGEMain.py} (\texttt{output/ige\_main.md}).} Within-race estimates are unweighted, which keeps the SEO
families; pooled estimates across races use the core longitudinal weights, without which the SEO would dominate the
pooled coefficient.\footnote{On the SRC sample alone the pooled IGE of sons' family income at 40 in 1980--2010 is
0.483, close to the 0.47 of \citet{palomino2018one}. Code: \texttt{5.CheckPMR.py}.} The main robustness check for
the racial comparison is the SRC sample alone, and each race weighted by its own sampling weights
(Table~\ref{tab:gaprobust}).

The children's adult years are split into six five-year bands, 25--29 to 50--55, and the 31--37 window of
\citet{chetty2020race} is added. Among sons, the bands contain 1,662 Black and 2,506 White sons at 25--29, 1,125 and
1,884 at 35--39, and 420 and 780 at 50--55. Older bands are observed only for earlier birth cohorts, but the
difference is smaller than it looks: the mean birth year is 1967 at 25--29 and 1962.5 at 45--49, and sons born 1951--64
make up between half and two thirds of every band.\footnote{\texttt{output/age\_profile\_robustness.md}, Section 0.
Code: \texttt{33.AgeProfileRobustness.py}.} Two features of the PSID matter for the oldest cohorts. Sons born
1951--56, a quarter to two fifths of each band, have parental income observed only at 17--19, because the survey starts
in 1968; the oldest bands therefore carry the noisiest parental measure, which attenuates the IGE where it is highest
and works against the age profile. And the panel loses members over time.

\paragraph{Attrition.} \citet{fitzgerald1998analysis} show that PSID attrition is selective on income. Among sons
observed at 25--29 and old enough to reach a later band, 58 percent of White and 40 percent of Black sons are observed
at 45--49 (Appendix Table~\ref{tab:attrition}). For White sons, staying rises with parental income, by 0.139$^{***}$ per log
point at 45--49 holding their own income at 25--29 fixed; for Black sons it does not (0.005). Selective attrition of
this kind could in principle create an opening gap, since the White sons who remain are those from richer families. It
does not: weighting each later band by the inverse of the predicted probability of staying changes the opening of
the gap by $-0.012$ (0.019) for earnings and $-0.014$ (0.016) for family income (Table~\ref{tab:gaprobust}). The
weights correct selection on parental income and on income at 25--29, not on unobserved wage growth. The NLSY79, whose
retention is far higher, gives the bound that a low-attrition panel provides: an opening of 0.088$^{**}$ (0.045) with
one year of parental income (Table~\ref{tab:threedata}).

The NLSY79 provides a second sample: respondents born 1960--64 who lived with a parent in 1979, with parental income
measured as the mean family income of the 1978--80 income years and adult outcomes followed to 55. The NLSY97 provides
a third: youths born 1980--84, with parental income reported by a parent for 1996, followed to their early forties. The
children of the NLSY79 women (NLSY79 Children and Young Adults) provide a fourth, with a multi-year parental measure:
the mother's family income averaged over the years in which the child was 13--19, and the child's wage and salary
income at each young-adult interview. It is a selected sample: children observed at 35--44 were born by the early
1980s to mothers who were then 23 or younger, and fewer than 300 sons of each race are observed at those
ages.\footnote{Code: \texttt{2.NLSY79.py}, \texttt{26.IGEByAge.py}, \texttt{29.NLSY97AgeProfile.py},
\texttt{10.CYAOutcomes.py} and \texttt{46.AgeProfileRound2.py}. The NLSY97 parental measure is a single year
for most youths; later years are reported by the youth and cover a selected, older subsample.}

The analysis of what the IGE contains (Section~\ref{sec:contains}) needs the father's own wage and hours and is
therefore restricted to sons whose father is the male head of their family at 16: 2,148 of 4,109 Black sons (52
percent) and 3,962 of 4,552 White sons (87 percent). On father-present sons the family-income IGE is 0.363 for Black
and 0.541 for White sons, against 0.268 and 0.473 on all sons.\footnote{Code: \texttt{24.IGEDecomposition.py}
(\texttt{output/ige\_decomposition.md}, Section 0).} The father's measures are means over his valid years when the son
is 12--18. Labor-market instability is measured by flags for displacement (plant closure, layoff or firing),
unemployment, reduced hours (below 1,500 while attached full time) and zero hours.

\subsection{Measures of persistence}

\paragraph{Log IGE by age band.} Within each age band $b$, I estimate
\begin{equation}\label{eq:ige}
y_{it} = \beta_{b}\, \bar y^{p}_{i} + f(A_i) + \delta_{a(i,t)} + \varepsilon_{it},
\end{equation}
where $y_{it}$ is the log outcome of child $i$ in year $t$ at an age in band $b$, $\bar y^{p}_{i}$ is log parental
income, $f(A_i)$ a quartic in the parent's age and $\delta_a$ single-year age dummies. All bands and both races are
estimated jointly in one stacked, fully interacted regression with standard errors clustered by 1968 family, so that
the change of $\beta_b$ across bands, and of the racial gap in it, can be tested. The summary statistic of the racial
divergence is the \emph{opening} of the gap: the mean White--Black gap over the bands from 35 minus the mean gap over
the bands below 35, on bands 25--49 unless stated (25--44 when the NLSY97 and the NLSY79 children are compared). All
racial differences in the paper are White minus Black. I also report the age-interacted design of \citet{lee2009trends}, which interacts parental income
with a quartic in the child's age and evaluates the IGE at every age.

\paragraph{Ranks and expected income.} Following \citet{chetty2014where} and \citet{chetty2020race}, the child's
band-mean family income is ranked within birth year and band, and parental income within the child's birth year,
among all children in the file; I also rank both generations within race. The dollar-weighted elasticity of
\citet{mitnik2020intergenerational} is estimated by Poisson pseudo-maximum likelihood on the income level, which keeps
years of zero earnings.

\paragraph{Reliability.} A single year of parental income measures permanent income with error. In the child
generation, the reliability ratio of one year of family income is 0.565 for Black and 0.585 for White parents in the
PSID (0.631 and 0.650 for the 12--17 mean), 0.556 and 0.552 in the NLSY79, and 0.492 and 0.430 in the
NLSY97.\footnote{Code: \texttt{4.Estimates.py}, following the logic of \citet{bolt2024intergenerational}; standard
errors 0.010--0.020.} Attenuation is common to all of the children's age bands, since parental income is measured
identically for every band: it scales the whole profile. The comparison across data sets is therefore made on the one
parental measure they share, one year, and then corrected race by race by each data set's own ratio, with the ratios'
sampling error carried into the standard errors.
